% !TeX root = ../main.tex
\section{Polymorphism}

STLC has a real problem: the identity function has to be written once per type. $\lambda x{:}\tyint.\,x$ and $\lambda x{:}\tybool.\,x$ are different terms with different types, even though they're doing the exact same thing. \textbf{Polymorphism} is the general name for language features that let one piece of code work at many types. Following the terminology Strachey and later Cardelli made standard, there are three classical kinds, and it's worth keeping them straight because they solve genuinely different problems:

\begin{itemize}
  \item \textbf{Parametric polymorphism}: one implementation, uniform behavior at every type (System F, Hindley-Milner generics).
  \item \textbf{Ad-hoc polymorphism}: one name, a \emph{different} implementation chosen per type (overloading, type classes, traits).
  \item \textbf{Subtype polymorphism}: one implementation, usable at any type that is a subtype of the one it expects. This is the subject of \S6 in full; we won't repeat it here.
\end{itemize}

\subsection{parametric polymorphism}

The fix for our identity-function problem is to abstract over the type itself, the same way $\lambda$ abstracts over a value. System F (also called the polymorphic lambda calculus) adds a universally quantified type $\forall \alpha.\, \tau$, together with a term-level construct that introduces one and one that eliminates it:
\[
  \tau ::= \dots \mid \forall \alpha.\, \tau
  \qquad
  e ::= \dots \mid \Lambda \alpha.\, e \mid e\,[\tau]
\]
$\Lambda \alpha.\, e$ (capital lambda) is \textbf{type abstraction}: it says ``$e$ works for any type $\alpha$, generic in $\alpha$''. $e\,[\tau]$ is \textbf{type application}: it instantiates that generic $e$ at a specific type $\tau$. The polymorphic identity function is now a single term with a single type:
\[
  \mathit{id} = \Lambda \alpha.\, \lambda x{:}\alpha.\, x \;:\; \forall \alpha.\, \alpha \to \alpha
\]
and $\mathit{id}\,[\tyint]\,3$ reduces (via type application, then ordinary $\beta$-reduction) to $3$, while $\mathit{id}\,[\tybool]\,\mathit{true}$ reduces to $\mathit{true}$ — the same underlying term, instantiated twice.

Writing out $\Lambda$ and $[\tau]$ explicitly everywhere gets tedious fast, which is exactly why Hindley-Milner (the type system behind ML, Haskell, and OCaml's core) exists: it's a restricted, decidable fragment of System F where quantifiers are only allowed at the very outside of a type ($\forall \alpha.\,\tau$ with $\tau$ quantifier-free — called a \emph{rank-1} or \emph{prenex} type) and instantiation is entirely implicit, inferred rather than written. That's the trade we made in \S4's Algorithm~W: give up some of System F's expressive power in exchange for type inference that always terminates and always finds a most general type.

\subsection{let-polymorphism vs.\ lambda-bound generalization}

Here's a wrinkle that trips people up the first time they see it. In ML-family languages,
\begin{lstlisting}[language=Haskell]
let id = \x -> x
in (id 3, id True)
\end{lstlisting}
typechecks fine — \lstinline|id| gets the polymorphic type $\forall\alpha.\,\alpha\to\alpha$ and each use instantiates it independently, once at $\tyint$ and once at $\tybool$. But
\begin{lstlisting}[language=Haskell]
(\id -> (id 3, id True)) (\x -> x)
\end{lstlisting}
does \emph{not} typecheck, even though it looks like it should do the same thing (apply the identity function to the same pair of uses). Why the difference?

The answer is that \textbf{generalization} — the step in Algorithm~W where a type gets its free variables universally quantified into a $\forall$ — only happens at \lstinline|let|, never at a $\lambda$-bound parameter. When we infer a type for the \lstinline|let|-bound \lstinline|id|, we typecheck its \emph{body} once, get a type like $\alpha \to \alpha$ with $\alpha$ a fresh, still-unconstrained variable, and generalize: since $\alpha$ doesn't appear anywhere in the surrounding context, we're free to quantify over it, giving $\mathit{id} : \forall \alpha.\,\alpha\to\alpha$, and every occurrence of \lstinline|id| in the \lstinline|let|-body gets to instantiate that quantifier fresh. A $\lambda$-bound parameter, by contrast, is unified with a \emph{single} type once and for all at every call site inside its scope — the type variable $\alpha$ standing for \lstinline|id|'s type gets unified with $\tyint\to\tyint$ at the first use (\lstinline|id 3|) and then, when the second use (\lstinline|id True|) tries to unify that same $\alpha$ with $\tybool\to\tybool$, unification fails.

This is exactly why \lstinline|let|-bindings are special-cased in HM (this is the "\textbf{let-polymorphism}" or "\textbf{$\mathrm{let}$-generalization}" rule) and it's also the reason behind the informal \textbf{value restriction} used in ML and OCaml: generalization is only sound when the right-hand side of the \lstinline|let| is a syntactic \emph{value} (a $\lambda$, a literal, a constructor application — something that doesn't perform an effect when evaluated). If the right-hand side were, say, a mutable reference constructor, generalizing its type could let you smuggle a value of one instantiation into a cell that later gets read back at an incompatible instantiation — unsound in the presence of mutable state. So: \lstinline|let x = e in ...| generalizes \lstinline|e|'s type only when \lstinline|e| is a value; otherwise it's treated monomorphically, just like a $\lambda$-bound parameter would be.

\subsection{ad-hoc polymorphism: overloading, type classes, traits}

Ad-hoc polymorphism looks similar from a distance — one name, many types — but it's solving a different problem: there's no single uniform implementation, just a family of \emph{different} ones, dispatched by type.

The crudest form is plain \textbf{overloading}: the same function name bound to unrelated implementations, resolved at compile time by argument type.
\begin{lstlisting}[language=Java]
int add(int a, int b) { return a + b; }
String add(String a, String b) { return a + b; }
\end{lstlisting}
There's no relationship enforced between the two \lstinline|add|s — the compiler just picks whichever overload matches the call site.

Haskell's \textbf{type classes} formalize this into something much more disciplined: a type class declares an interface, and instances supply implementations per type, with the compiler inferring \emph{which} instance to use from the call site's inferred type.
\begin{lstlisting}[language=Haskell]
class Eq a where
  (==) :: a -> a -> Bool

instance Eq Int where
  x == y = intEq x y

instance Eq Bool where
  True  == True  = True
  False == False = True
  _     == _     = False

-- (==) works for any type with an Eq instance:
allEqual :: Eq a => [a] -> Bool
allEqual xs = and (zipWith (==) xs (tail xs))
\end{lstlisting}
\lstinline|allEqual| is itself parametrically polymorphic in \lstinline|a| \emph{constrained} by an ad-hoc-polymorphic operation — this combination (parametric polymorphism plus a typeclass constraint) is sometimes called \textbf{bounded polymorphism}, and it's genuinely the best of both worlds: uniform code, per-type behavior, resolved automatically.

Rust's \textbf{traits} are the same idea under a different name, with dispatch resolved either statically (monomorphization, one specialized copy of the generic function per concrete type, the default) or dynamically (\lstinline|dyn Trait|, a vtable, chosen with an explicit annotation):
\begin{lstlisting}
trait Speak {
    fn speak(&self) -> String;
}

impl Speak for Dog {
    fn speak(&self) -> String { "Woof".to_string() }
}

impl Speak for Cat {
    fn speak(&self) -> String { "Meow".to_string() }
}

fn greet<T: Speak>(x: &T) -> String {
    format!("Hello, {}", x.speak())
}
\end{lstlisting}
\lstinline|greet| is generic over any \lstinline|T| implementing \lstinline|Speak| — parametric in \lstinline|T|, ad-hoc in which \lstinline|speak| gets called.

\subsection{higher-rank polymorphism}

Plain Hindley-Milner only quantifies types at the very top (rank-1 / prenex): $\forall\alpha.\,\tau$ where $\tau$ contains no further $\forall$. This is a real restriction, and it bites the moment you want to pass a polymorphic function \emph{as an argument} to another function and have it stay polymorphic once inside.

\begin{example}[where plain HM inference gets stuck]
Suppose you want to write a function that applies a polymorphic function to two different types of argument:
\begin{lstlisting}[language=Haskell]
apply2 :: (forall a. a -> a) -> (Int, Bool)
apply2 f = (f 3, f True)
\end{lstlisting}
The parameter \lstinline|f| itself needs type $\forall\alpha.\,\alpha\to\alpha$ — a quantifier appearing \emph{inside} the argument position of \lstinline|apply2|'s own type, not just out front. That's a \textbf{rank-2} type: $(\forall\alpha.\,\alpha\to\alpha) \to (\tyint \times \tybool)$. Plain Algorithm~W cannot infer this. If you write \lstinline|apply2 (\x -> x)| without the annotation, HM will unify the type of \lstinline|x -> x| against a single monomorphic use as soon as it sees \lstinline|f 3|, fixing \lstinline|f|'s type at $\tyint\to\tyint$ — and then \lstinline|f True| fails to typecheck, for exactly the same reason the $\lambda$-bound \lstinline|id| failed above. The fix in Haskell is to write the rank-2 signature explicitly (as we did above): once the type is given, \emph{checking} against it (rather than inferring it from scratch) works fine — this is one of the places where full type inference is undecidable but type checking against a supplied signature is not.
\end{example}

This is System F$_\omega$ territory: dropping the rank-1 restriction and allowing $\forall$ anywhere in a type gives you the full expressive power of System F, at the cost of losing decidable, fully-automatic type inference in general. Real languages compromise: Haskell lets you opt in to higher-rank types via explicit annotations (the \texttt{RankNTypes} extension) while keeping ordinary rank-1 inference as the default, unannotated case.

\subsection{Exercises}
\begin{itemize}
  \item Write the System F term and type for a polymorphic \lstinline|const| function ($\mathit{const}\ x\ y = x$), using explicit $\Lambda$ and $[\tau]$. What is $\mathit{const}\,[\tyint]\,[\tybool]\;3\;\mathit{true}$'s normal form?
  \item Explain, in your own words, why \lstinline|let f = (\x -> x) in (f 1, f True)| typechecks but \lstinline|(\f -> (f 1, f True)) (\x -> x)| does not, referring to when generalization happens in Algorithm~W.
  \item The value restriction blocks generalizing the type of \lstinline|let r = ref [] in ...| in ML (where \lstinline|ref| allocates a mutable cell). Construct a small example showing what would go wrong, type-safety-wise, if \lstinline|r|'s type \emph{were} generalized to $\forall\alpha.\,\mathbf{ref}\,(\mathbf{List}\,\alpha)$.
  \item Implement the same ad-hoc-polymorphic behavior (say, a \lstinline|describe| function for at least two unrelated types) three ways: Java-style overloading, a Haskell type class, and a Rust trait. Which one lets you add a new case \emph{without} touching the original definition?
  \item Why can't Algorithm~W infer a type for \lstinline|apply2 (\x -> x)| in the example above without an explicit rank-2 annotation? Point to the specific unification step that goes wrong.
\end{itemize}
